\DocumentMetadata{
  pdfversion=2.0,
  pdfstandard=ua-2,
  lang=en-US,
  tagging=on,
  tagging-setup={math/setup=mathml-SE,tabsorder=structure}
}
\documentclass[11pt,letterpaper]{article}
\usepackage[margin=0.68in,includefoot]{geometry}
\usepackage{newcomputermodern}
\usepackage{amsmath,amscd,mathtools}
\usepackage{tikz,adjustbox}
\providecommand{\square}{\mdlgwhtsquare}
\usepackage[hidelinks]{hyperref}
\hypersetup{
  pdftitle={Numerical Analysis PhD Exam, August 2016, Part 1},
  pdfauthor={University of Florida Department of Mathematics},
  pdfsubject={Graduate mathematics examination},
  pdfdisplaydoctitle=true,
  bookmarksopen=true
}
\setcounter{secnumdepth}{0}
\setlength{\parindent}{0pt}
\setlength{\parskip}{0.28em}
\setlength{\emergencystretch}{3em}
\setlength{\leftmargini}{2.15em}
\setlength{\leftmarginii}{2.65em}
\setlength{\leftmarginiii}{3em}
\pagestyle{empty}
\raggedbottom
\allowdisplaybreaks
\NewTaggingSocketPlug{block/list/label}{examproblem}{%
  \tagstructbegin{tag=Lbl}\tagstructend%
  \tagstructbegin{tag=\LBody}%
  \tagstructbegin{tag=\ProblemHeadingTag,title-o={\ProblemHeadingText}}%
  \tagmcbegin{tag=\ProblemHeadingTag,actualtext={\ProblemHeadingText}}%
  #2%
  \tagmcend%
  \tagstructend%
}
\newenvironment{examproblems}[2]{%
  \def\ProblemHeadingTag{#2}%
  \AssignTaggingSocketPlug{block/list/label}{examproblem}%
  \begin{enumerate}%
  \setcounter{enumi}{#1}%
  \renewcommand{\labelenumi}{\textbf{\theenumi.}}%
  \setlength{\topsep}{0.35em}%
  \setlength{\partopsep}{0pt}%
  \setlength{\itemsep}{0.48em}%
  \setlength{\parsep}{0.18em}%
}{\end{enumerate}}
\newenvironment{examsubparts}{%
  \AssignTaggingSocketPlug{block/list/label}{default}%
  \begin{enumerate}%
  \setlength{\topsep}{0.18em}%
  \setlength{\partopsep}{0pt}%
  \setlength{\itemsep}{0.16em}%
  \setlength{\parsep}{0.1em}%
}{\end{enumerate}}
\newenvironment{examinstructions}{%
  \par\begingroup\itshape
}{%
  \par\endgroup\vspace{0.18em}
}
\makeatletter
\renewcommand\subsection{\@startsection{subsection}{2}{0pt}%
  {-1.25ex plus -.25ex minus -.1ex}%
  {0.55ex plus .1ex}%
  {\normalfont\large\bfseries}}
\renewcommand{\@maketitle}{%
  \newpage\null\vskip -1em%
  {\footnotesize\bfseries
    \href{https://www.ufl.edu/}{University of Florida}, \href{https://math.ufl.edu/}{Department of Mathematics}\par}%
  \vskip -0.5em%
  \tagstructbegin{tag=H1,title={Numerical Analysis PhD Exam, August 2016, Part 1}}%
  \tagpdfparaOff%
  \tagmcbegin{tag=H1}%
    {\LARGE\bfseries\href{https://gma.math.ufl.edu/exams/numerical-analysis-phd/}{Numerical Analysis PhD Exam}, August 2016, Part 1\par}%
  \tagmcend%
  \tagpdfparaOn%
  \tagstructend%
  \vskip 0.25em%
  \tagmcbegin{artifact}%
    \hrule height 1pt%
  \tagmcend%
  \vskip 1em%
}
\makeatother
\title{Numerical Analysis PhD Exam, August 2016, Part 1}
\author{}
\date{}
\begin{document}
\maketitle
\thispagestyle{empty}
\begin{examinstructions}
Do 4 (four) problems.
\end{examinstructions}
\begin{examproblems}{0}{H2}
\def\ProblemHeadingText{Problem 1.}
\item
Define a normal matrix and prove that the following are equivalent.
\begin{examsubparts}
\item[\textnormal{(a)}]
\(A\) is normal.
\item[\textnormal{(b)}]
\(A\) is unitarily diagonalizable.
\item[\textnormal{(c)}]
\(\|A\|_F=\left(\sum |\lambda_i|^2\right)^{1/2}\), where \(\{\lambda_i\}\) are the eigenvalues of~\(A\) counted with multiplicity.
\end{examsubparts}
\def\ProblemHeadingText{Problem 2.}
\item
Let \(\kappa_2(A)\) be the two-norm condition number of the square, non-singular~\(A\).
\begin{examsubparts}
\item[\textnormal{(a)}]
Prove that
\[
\kappa_2(A)=\frac{\sigma_1}{\sigma_m},
\]
 where \(\sigma_1\) and \(\sigma_m\) are the largest and smallest singular values of~\(A\), respectively.
\item[\textnormal{(b)}]
Prove or disprove: If \(A=QBQ^*\) with~\(Q\) unitary, then \(\kappa_2(A)=\kappa_2(B)\).
\item[\textnormal{(c)}]
Prove or disprove: If \(A=CBC^{-1}\), then \(\kappa_2(A)=\kappa_2(B)\).
\end{examsubparts}
\def\ProblemHeadingText{Problem 3.}
\item
Assume \(A\in\mathbb{R}^{m,n}\) with \(m\geq n\), \(\operatorname{rank}(A)=n\), and \(b\in\mathbb{R}^m\).
\begin{examsubparts}
\item[\textnormal{(a)}]
Define the least squares solution to \(Ax=b\).
\item[\textnormal{(b)}]
Derive the normal equations for the least squares problem.
\item[\textnormal{(c)}]
Prove that \(A^T A\) is invertible.
\item[\textnormal{(d)}]
Prove that the unique solution to the least squares problem is \((A^T A)^{-1}A^T b\).
\item[\textnormal{(e)}]
Describe how to solve the least squares problem using the QR decomposition of~\(A\).
\end{examsubparts}
\def\ProblemHeadingText{Problem 4.}
\item
This problem has two parts.
\begin{examsubparts}
\item[\textnormal{(a)}]
Prove that~\(P\) is an orthogonal projector if and only if it is Hermitian.
\item[\textnormal{(b)}]
Let \(\{q_1,q_2,\ldots,q_n\}\) be an orthonormal subset of \(\mathbb{C}^m\). Show that
\[
P=\sum_{i=1}^n q_iq_i^*
\]
 is an orthogonal projector with range equal to the span of \(\{q_1,q_2,\ldots,q_n\}\).
\end{examsubparts}
\def\ProblemHeadingText{Problem 5.}
\item
Assume \(A\in\mathbb{R}^{m,m}\).
\begin{examsubparts}
\item[\textnormal{(a)}]
Prove that \(\langle x,y\rangle_A=x^T Ay\) is an inner product on \(\mathbb{R}^m\) if and only if~\(A\) is symmetric and positive definite.
\item[\textnormal{(b)}]
Assume now that~\(A\) is symmetric and positive definite. If \(x_*\) is the solution to \(Ax=b\) and \(\{p_1,\ldots,p_m\}\) is an orthonormal basis for \(\mathbb{R}^m\) with respect to \(\langle\, ,\,\rangle_A\) and \(x_*=\sum c_i p_i\), give a formula for the \(c_i\).
\end{examsubparts}
\end{examproblems}
\end{document}
